\begin{abstract}
    Cooperative multi-agent reinforcement learning (MARL) with parameter sharing performs well on multi-warehouse inventory control under symmetric conditions. However, real distribution networks typically exhibit heterogeneity in demand, product mix, and shipment cost. We introduce a single-echelon benchmark with an order-level demand model and four throughput-matched heterogeneous environments to investigate whether a parameter-shared policy degrades under heterogeneity and if so, what interventions recover lost performance. We find that although per-environment re-tuning brings the policy close to a simulation-optimized heuristic, a substantial performance gap remains compared to a centralized full-information policy and we trace this residual gap to the two distinct failure modes of misallocating product service levels and failing to scale order quantities to regional demand. To close the gap, we test five modifications to the shared-policy pipeline and find that the useful gains come from cheap, targeted controls rather than added architecture. Specifically, Per-warehouse normalization, a static weight feature, and per-SKU credit assignment each recover a modest, axis-specific part of the gap, while the per-SKU attention architecture we proposed as the main remedy does not     pay off at this scale and is matched by a flat network of equal capacity. Our practical takeaway for heterogeneous settings is to diagnose the active heterogeneity dimension and apply the cheapest control that targets it. Per-SKU attention is worth revisiting only at larger scale where its constant parameter count begins to matter.
    \end{abstract}